\chapter{Orthogonal Projection and Least Squares}
\label{ch:projection-least-squares}

\section{The Best Approximation Problem}

The equation $Ax=b$ asks whether $b$ is in the column space of $A$. If it is
not, there is no exact solution. Least squares asks a different question:

\[
  \text{find }x\text{ so that }Ax\text{ is as close as possible to }b.
\]

This is not a trick for avoiding linear algebra. It is linear algebra. We are
looking for the point in $\Img(A)$ closest to $b$.

\begin{definition}[Least-squares problem]
For $A \in \R^{m \times n}$ and $b \in \R^m$, a \emph{least-squares solution}
is a vector $\hat{x}\in\R^n$ minimizing
\[
  \norm{Ax-b}.
\]
The vector $\hat{b}=A\hat{x}$ is the least-squares approximation to $b$ from
the column space of $A$, and
\[
  r=b-A\hat{x}
\]
is the residual.
\end{definition}

The unknown $\hat{x}$ may not be unique. The fitted vector $\hat{b}$ is unique:
there is only one closest point to $b$ in a finite-dimensional subspace.

\section{Projection Onto A Subspace}

\begin{definition}[Orthogonal complement]
If $U$ is a subspace of an inner product space $V$, its orthogonal complement is
\[
  U^\perp=\{v\in V:\inner{v}{u}=0\text{ for every }u\in U\}.
\]
\end{definition}

\begin{proposition}
If $U$ is a subspace of a finite-dimensional inner product space $V$, then
\[
  V = U \oplus U^\perp.
\]
Thus every $v\in V$ can be written uniquely as
\[
  v=u+w,\qquad u\in U,\quad w\in U^\perp.
\]
\end{proposition}

\begin{proof}
Choose an orthonormal basis $q_1,\ldots,q_k$ for $U$. For any $v\in V$, set
\[
  u=\sum_{j=1}^k \inner{v}{q_j}q_j.
\]
Then $u\in U$. For each $i$,
\[
  \inner{v-u}{q_i}
  =
  \inner{v}{q_i}
  -
  \sum_{j=1}^k \inner{v}{q_j}\inner{q_j}{q_i}
  =
  \inner{v}{q_i}-\inner{v}{q_i}
  =
  0.
\]
Since $q_1,\ldots,q_k$ span $U$, we have $v-u\in U^\perp$. So
$v=u+(v-u)$.

For uniqueness, if $z\in U\cap U^\perp$, then $\inner{z}{z}=0$, so $z=0$.
Thus the sum is direct.
\end{proof}

\begin{definition}[Orthogonal projection onto a subspace]
The vector $u$ in the decomposition $v=u+w$, with $u\in U$ and
$w\in U^\perp$, is the \emph{orthogonal projection} of $v$ onto $U$. We write
\[
  u=\Proj_U(v).
\]
\end{definition}

\begin{proposition}[Best approximation property]
Let $U$ be a subspace of a finite-dimensional inner product space $V$. For any
$v\in V$, $\Proj_U(v)$ is the unique vector in $U$ minimizing
\[
  \norm{v-u},\qquad u\in U.
\]
\end{proposition}

\begin{proof}
Write $v=\hat{u}+r$ with $\hat{u}\in U$ and $r\in U^\perp$. For any other
$u\in U$,
\[
  v-u = r + (\hat{u}-u).
\]
Here $r\perp(\hat{u}-u)$ because $\hat{u}-u\in U$. By the Pythagorean theorem,
\[
  \norm{v-u}^2
  =
  \norm{r}^2+\norm{\hat{u}-u}^2
  \geq
  \norm{r}^2
  =
  \norm{v-\hat{u}}^2.
\]
Equality forces $\hat{u}=u$.
\end{proof}

\section{Least Squares As Projection}

Let $U=\Img(A)\subseteq\R^m$. Least squares is the projection problem
\[
  \hat{b}=\Proj_{\Img(A)}(b).
\]
If $\hat{b}=A\hat{x}$, then the residual
\[
  r=b-A\hat{x}
\]
is orthogonal to the whole column space of $A$.

\begin{theorem}[Normal equations]
A vector $\hat{x}$ is a least-squares solution to $Ax=b$ if and only if
\[
  A^T(A\hat{x}-b)=0.
\]
Equivalently,
\[
  A^TA\hat{x}=A^Tb.
\]
\end{theorem}

\begin{proof}
The vector $A\hat{x}$ is the closest point to $b$ in $\Img(A)$ exactly when
the residual $b-A\hat{x}$ is orthogonal to $\Img(A)$. The columns of $A$ span
$\Img(A)$, so this is equivalent to
\[
  a_j^T(b-A\hat{x})=0
  \quad\text{for every column }a_j\text{ of }A.
\]
Stacking these conditions gives
\[
  A^T(b-A\hat{x})=0,
\]
which is the same as $A^TA\hat{x}=A^Tb$.
\end{proof}

\begin{warningbox}
The normal equations are mathematically central, but they are not always the
best numerical method. Forming $A^TA$ can amplify conditioning problems. The
geometry belongs here; the numerically safer QR route belongs to
Chapter~\ref{ch:regression-qr-stability}.
\end{warningbox}

\section{When Is The Least-Squares Solution Unique?}

The fitted vector $\hat{b}=\Proj_{\Img(A)}(b)$ is always unique. The coefficient
vector $\hat{x}$ is unique exactly when the columns of $A$ are linearly
independent.

\begin{proposition}
The following are equivalent for $A\in\R^{m\times n}$:
\begin{enumerate}[label=(\alph*)]
  \item the least-squares solution is unique for every $b\in\R^m$;
  \item $A^TA$ is invertible;
  \item $\Ker(A)=\{0\}$;
  \item the columns of $A$ are linearly independent.
\end{enumerate}
\end{proposition}

\begin{proof}
The key identity is
\[
  \Ker(A^TA)=\Ker(A).
\]
Indeed, if $Av=0$, then $A^TAv=0$. Conversely, if $A^TAv=0$, then
\[
  0=v^TA^TAv=\norm{Av}^2,
\]
so $Av=0$. Thus $A^TA$ is invertible exactly when $\Ker(A)=\{0\}$, which is
exactly column independence. In that case the normal equations have exactly one
solution for every $b$.
\end{proof}

\section{A Small Computation By Hand}

Fit a line $y=a+bt$ to the points
\[
  (0,1),\qquad (1,2),\qquad (2,2).
\]
The matrix equation is
\[
  \begin{bmatrix}
  1&0\\
  1&1\\
  1&2
  \end{bmatrix}
  \begin{bmatrix}a\\ b\end{bmatrix}
  =
  \begin{bmatrix}1\\2\\2\end{bmatrix}.
\]
The system is overdetermined. The normal equations are
\[
  \begin{bmatrix}
  3&3\\
  3&5
  \end{bmatrix}
  \begin{bmatrix}a\\b\end{bmatrix}
  =
  \begin{bmatrix}5\\6\end{bmatrix}.
\]
Solving gives
\[
  \hat{a}=\frac{7}{6},
  \qquad
  \hat{b}=\frac{1}{2}.
\]
The fitted line is
\[
  \hat{y}=\frac{7}{6}+\frac{1}{2}t.
\]

\section{The Four Fundamental Subspaces}

Orthogonality also organizes the row space, column space, null space, and left
null space of a matrix.

\begin{definition}[Four fundamental subspaces]
For $A\in\R^{m\times n}$, the four fundamental subspaces are
\[
  \col(A)=\Img(A)\subseteq\R^m,
  \qquad
  \Ker(A)\subseteq\R^n,
\]
\[
  \row(A)=\col(A^T)\subseteq\R^n,
  \qquad
  \Ker(A^T)\subseteq\R^m.
\]
\end{definition}

The row space and null space live in $\R^n$. The column space and left null
space live in $\R^m$. Keeping the ambient spaces straight prevents many
mistakes.

\begin{theorem}[Fundamental orthogonal decompositions]
For $A\in\R^{m\times n}$,
\[
  \R^n = \row(A)\oplus \Ker(A),
  \qquad
  \R^m = \col(A)\oplus \Ker(A^T).
\]
\end{theorem}

\begin{proof}
The equation $Ax=0$ means every row of $A$ has dot product zero with $x$.
Therefore
\[
  \Ker(A)=\row(A)^\perp.
\]
Since $\R^n=U\oplus U^\perp$ for any subspace $U\subseteq\R^n$, taking
$U=\row(A)$ gives the first decomposition.

Applying the same argument to $A^T$ gives
\[
  \Ker(A^T)=\row(A^T)^\perp=\col(A)^\perp,
\]
so $\R^m=\col(A)\oplus\Ker(A^T)$.
\end{proof}

\begin{computebox}
These decompositions explain solvability. A vector $b\in\R^m$ splits uniquely
as
\[
  b=b_{\col}+b_{\perp},
  \qquad
  b_{\col}\in\col(A),\quad b_{\perp}\in\Ker(A^T).
\]
The equation $Ax=b$ is solvable exactly when $b_{\perp}=0$. Least squares
keeps $b_{\col}$ and discards the orthogonal obstruction $b_{\perp}$.
\end{computebox}

\section{Projection Matrices}

When $A$ has linearly independent columns, the projection of $b$ onto
$\col(A)$ is
\[
  \hat{b}=A(A^TA)^{-1}A^Tb.
\]
The matrix
\[
  H=A(A^TA)^{-1}A^T
\]
is called the projection matrix or hat matrix. It sends $b$ to $\hat{b}$.

\begin{proposition}
If $A$ has linearly independent columns, then
\[
  H^2=H,\qquad H^T=H.
\]
\end{proposition}

\begin{proof}
For idempotence,
\[
  H^2
  =
  A(A^TA)^{-1}A^T A(A^TA)^{-1}A^T
  =
  A(A^TA)^{-1}A^T
  =
  H.
\]
For symmetry, use $(A^TA)^T=A^TA$, so $(A^TA)^{-1}$ is also symmetric:
\[
  H^T
  =
  \left(A(A^TA)^{-1}A^T\right)^T
  =
  A(A^TA)^{-1}A^T.
\]
\end{proof}

\section{Python Thread}

\begin{lstlisting}
import numpy as np

A = np.array([[1.0, 0.0],
              [1.0, 1.0],
              [1.0, 2.0]])
b = np.array([1.0, 2.0, 2.0])

# Least squares by normal equations
x_hat = np.linalg.solve(A.T @ A, A.T @ b)
b_hat = A @ x_hat
r = b - b_hat

print("x_hat:", x_hat)
print("residual:", r)
print("A.T @ residual:", A.T @ r)

# Projection matrix
H = A @ np.linalg.solve(A.T @ A, A.T)
print("H^2 = H:", np.allclose(H @ H, H))
print("H.T = H:", np.allclose(H.T, H))
print("H @ b equals b_hat:", np.allclose(H @ b, b_hat))
\end{lstlisting}

\begin{labbox}
Lab 7 asks you to compute the four fundamental subspaces numerically and check
the two orthogonal decompositions. Keep the dimensions visible: row and null
spaces live in the domain; column and left-null spaces live in the codomain.
\end{labbox}

\section*{Chapter Summary}
\addcontentsline{toc}{section}{Chapter Summary}

\begin{itemize}
  \item Least squares is projection onto the column space.
  \item The residual of a least-squares approximation is orthogonal to the
        column space.
  \item The normal equations are $A^TA\hat{x}=A^Tb$.
  \item The fitted vector is unique; the coefficient vector is unique when
        the columns of $A$ are independent.
  \item The row/null and column/left-null decompositions explain the geometry
        of solving $Ax=b$.
\end{itemize}

\section*{Exercises}
\addcontentsline{toc}{section}{Exercises}

\subsection*{A. Theory}

\begin{exercise}
\exerciseid{8.A1}
Prove that if $\hat{x}$ minimizes $\norm{Ax-b}$, then the residual
$b-A\hat{x}$ is orthogonal to every column of $A$.
\end{exercise}

\begin{exercise}
\exerciseid{8.A2}
Prove that $\Ker(A^TA)=\Ker(A)$ for every real matrix $A$.
\end{exercise}

\subsection*{B. By Hand}

\begin{exercise}
\exerciseid{8.B1}
Project $b=(2,1,0)$ onto the line spanned by $a=(1,1,1)$ and compute the
residual. Check that the residual is orthogonal to $a$.
\end{exercise}

\begin{exercise}
\exerciseid{8.B2}
For
\[
  A=\begin{bmatrix}1&0\\1&1\\1&2\end{bmatrix},
  \qquad
  b=\begin{bmatrix}1\\2\\2\end{bmatrix},
\]
write down the normal equations and solve them by hand.
\end{exercise}

\subsection*{C. Python / NumPy}

\begin{exercise}
\exerciseid{8.C1}
For a small overdetermined matrix $A$, use \texttt{np.linalg.lstsq} and verify
numerically that \texttt{A.T @ (b - A @ x)} is close to zero.
\end{exercise}

\begin{exercise}
\exerciseid{8.C2}
For a full-column-rank matrix $A$, build
\[
  H=A(A^TA)^{-1}A^T
\]
using \texttt{np.linalg.solve}. Check numerically that $H^2=H$ and $H^T=H$.
\end{exercise}

\subsection*{D. Challenge}

\begin{exercise}
\exerciseid{8.D1}
Give an example of a matrix $A$ and a vector $b$ for which $Ax=b$ has no
solution. Then decompose $b$ as $b_{\col}+b_\perp$, where
$b_{\col}\in\col(A)$ and $b_\perp\in\Ker(A^T)$.
\end{exercise}

